% ==================================================================
% id:       comb-001
% topic:    קומבינטוריקה
% subtopic: עקרון הכפל וסידורים
% solution: full
% ==================================================================
פרופסור אקס רוצה לסדר על המדף במשרדו החדש 10 ספרים \underline{שונים}: 4 ספרי היסטוריה, 3 ספרי מדע בדיוני, 2 ספרי רפואה וספר סטטיסטיקה אחד.
\begin{enumerate}
\item בכמה דרכים הוא יכול לעשות זאת?
\begin{solbox}
אלה 10 ספרים שונים, ולכן יש $\ans{10!}$ דרכים לסדר אותם על המדף.
\end{solbox}
\item בכמה דרכים הוא יכול לעשות זאת, אם הוא רוצה שספר הסטטיסטיקה יהיה באחד מקצוות המדף?
\begin{solbox}
מספר הסידורים שבהם ספר הסטטיסטיקה בקצה הימני הוא $9!$, כי נשאר לסדר את 9 הספרים האחרים. באותו אופן, יש $9!$ סידורים שבהם הוא בקצה השמאלי. בסך הכול: $\ans{2\cdot9!}$.
\end{solbox}
\item בכמה דרכים הוא יכול לעשות זאת, אם הוא רוצה שכל הספרים באותו נושא יהיו צמודים זה לזה?
\begin{solbox}
נחשוב על כל הספרים באותו נושא כעל "בלוק" אחד. יש 4 בלוקים לסדר בשורה, ואת זה אפשר לעשות ב-$4!$ דרכים. לכל סידור כזה: את ספרי ההיסטוריה אפשר לסדר בתוך הבלוק ב-$4!$ דרכים, את ספרי המדע הבדיוני ב-$3!$, את ספרי הרפואה ב-$2!$ ואת ספר הסטטיסטיקה ב-$1!$ דרכים. לפי עקרון הכפל:
\[
\ans{4!\cdot4!\cdot3!\cdot2!\cdot1!}.
\]
\end{solbox}
\item אחייניתו של פרופסור אקס מגיעה לביקור, וכמתנה ליום הולדתה השביעי הוא רוצה לתת לה שלושה מתוך עשרת הספרים. כמה מתנות שונות הוא יכול לתת לה?
\begin{solbox}
מתנה נקבעת לפי הרכב הספרים, בלי חשיבות לסדר ובלי חזרות. לכן יש $\ans{\binom{10}3}$ מתנות שונות.
\end{solbox}
\item בכל יום ראשון הפרופסור בוחר שני ספרים ושם אותם בתיבת הדואר של אגודת הסטודנטים. כפרופסור מסודר, הוא מנהל רשימה שבה הוא רושם את שמות שני הספרים שבחר בכל שבוע.
\begin{enumerate}
\item בהנחה שהסטודנטים לא מחזירים לפרופסור את הספרים, כמה רשימות שונות הוא יכול לנהל עד שייגמרו לו הספרים (כלומר בתום 5 שבועות)?
\item בהנחה שהסטודנטים מחזירים את הספרים עד סוף כל שבוע, כמה רשימות שונות יכולות להיות לפרופסור בתום 5 שבועות?
\end{enumerate}
\begin{solbox}
\begin{enumerate}
\item בכל שבוע בוחרים זוג ספרים, בלי חשיבות לסדר ובלי חזרות. בשבוע הראשון יש $\binom{10}2$ אפשרויות, בשני $\binom82$, בשלישי $\binom62$, ברביעי $\binom42$ ובחמישי $\binom22=1$. לפי עקרון הכפל:
\[
\ans{\binom{10}2\binom82\binom62\binom42\binom22}.
\]
\item כאן בכל שבוע מספר האפשרויות נשאר $\binom{10}2$, כי הספרים חוזרים לפני הבחירה הבאה. לכן יש $\ans{\binom{10}2^5}$ רשימות.
\end{enumerate}
\end{solbox}
\end{enumerate}
